\documentclass[10pt,a4paper]{amsart}
\usepackage[margin=19mm]{geometry}
\usepackage{amsmath,amssymb,mathtools}
\usepackage[scr=rsfs,cal=euler]{mathalfa}
\usepackage{xfrac}
\usepackage[unicode,hidelinks,bookmarks=false]{hyperref}
\newcommand{\ie}{i.e.\ }
\newcommand{\cf}{cf.\ }
\newcommand{\resp}{respectively\ }
\newcommand{\Subsection}{\subsection}
\providecommand{\nobreakdash}{\nobreak\hbox{-}}
\setlength{\emergencystretch}{2em}
\multlinegap0pt
\newcommand{\C}{\mathbb{C}}
\usepackage{mathtools}
\newtheorem{theorem}{Theorem}
\newtheorem{lemma}{Lemma}
\newtheorem{corollary}{Corollary}
\newtheorem{proposition}{Proposition}
\title{Convexity of the Berezin range of composition operators induced by automorphisms of the Unit Ball in $\C^{n}$}
\author{Timothy G. Clos, Trevor C. Pentzien,, Tomas Miguel Rodriguez}
\date{Source: arXiv:2609.10361v1 (CC BY 4.0)}
\begin{document}
\maketitle

\noindent\emph{Source and licence.} Convexity of the Berezin range of composition operators induced by automorphisms of the Unit Ball in $\C^{n}$. Version arXiv:2609.10361v1 (CC BY 4.0), distributed by arXiv under the Creative Commons Attribution 4.0 International licence, http://creativecommons.org/licenses/by/4.0/, as recorded in the arXiv licence metadata for this exact version and shown on the abstract page https://arxiv.org/abs/2609.10361v1. Full licence notice: \texttt{licenses/arxiv-2609.10361-CC-BY-4.0.txt}.\par\smallskip

\noindent\textit{Editorial scope.} Counted unit: the article's theorem MainResultOneDimension (source0.tex lines 118--120). The article's complete original proof of it is delivered with it (source0.tex lines 281--365, one \texttt{proof} environment of 85 lines, which the article places away from the statement). No mathematics is written, repaired or re-derived: the statement and the proof are the article's own lines, in the article's own notation, and no retained line is substituted or re-flowed.  Retained uncounted as the source-local context the proof needs: lines 160--169; lines 257--259; lines 265--267.  Nothing was omitted from any retained span, so the package carries no omission record.  No retained line contains a citation, so the wrapper carries no bibliography and no bibliography entry.  No retained line contains a figure environment, an image-inclusion command or drawing code, so no image is delivered and the package declares no asset dependency.  Editorial formatting only, in addition: an AMS \texttt{amsart} preamble that restates the article's own declarations and macros used by the retained lines, namely the environments \texttt{theorem}, \texttt{lemma}, \texttt{corollary}, \texttt{proposition} on separate counters, together with the standard packages those lines need.  The retained environments print in this document as Theorem 1, Lemma 1, Corollary 1, Proposition 1 in order of appearance, whereas the article prints the same items with the numbering of its own full document.  The complete native source of the article as distributed in the arXiv e-print for this exact version is bundled unchanged as \texttt{sources/209903/source0.tex}.
\begin{lemma}\label{ReImBerezinTransform}
The real and imaginary parts of $\widetilde{C}_{\varphi_{a,\xi}}$ are of the form,
\[\mathrm{Re}(\widetilde{C}_{\varphi_{a,\xi}})(z)=\left[\dfrac{1-|z|^2}{A(z)^2+B(z)^2}\right]^2\left(C(z)[A(z)^2-B(z)^2]+4A(z)B(z)D(z)\right),\]
\[\mathrm{Im}(\widetilde{C}_{\varphi_{a,\xi}})(z)=2 \left[\dfrac{1-|z|^2}{A(z)^2+B(z)^2}\right]^2(D(z)[A(z)^2-B(z)^2]-A(z)B(z)C(z)),\]
where
\[A(z)= -\mathrm{Re}(\overline{a}z)(1+\mathrm{Re}(\xi))+ 1 + \mathrm{Re}(\xi)|z|^2-\mathrm{Im}(\xi)\mathrm{Im}(\overline{a}z),\]
\[B(z)=-\mathrm{Im}(\overline{a}z)(1-\mathrm{Re}(\xi))+\mathrm{Im}(\xi)|z|^2-\mathrm{Im}(\xi)\mathrm{Re}(\overline{a}z),\]
\[C(z)=(1-\mathrm{Re}(\overline{a}z))^2-\mathrm{Im}(\overline{a}z)^2,\]
\[D(z)=-(1-\mathrm{Re}(\overline{a}z))\mathrm{Im}(\overline{a}z).\]
\end{lemma}    

\begin{corollary}\label{propconj}
Let $\xi \in \mathbb{T}$.  $B(C_{\varphi_{a,\xi}})$ is closed under complex conjugation if and only if $\mathrm{Im}(\xi) = 0$.
\end{corollary}

\begin{proposition}\label{auto_on_disk}
    Let $\tau$ be a rotation of $\mathbb{D}$.  Then the Berezin range $B(C_{\varphi_{\tau(a),\xi}})=B(C_{\varphi_{a,\xi}})$ as sets. 
\end{proposition}

\begin{theorem}\label{MainResultOneDimension}
The Berezin range of $C_{\varphi_{a,1}}$ on $A^2(\mathbb{D})$ is convex if and only if $0\leq|a| \leq \frac{\sqrt{2}}{2}$. 
\end{theorem}

\begin{proof}[Proof of Theorem \ref{MainResultOneDimension}]
By Proposition \ref{auto_on_disk} without loss of generality we can assume that $a$ is a nonnegative real number.
Suppose $B(C_{\varphi_{a,1}})$ is convex, we want to show that $0\leq a \leq \frac{\sqrt{2}}{2}$.
Assume the contrary. Using Lemma \ref{ReImBerezinTransform} and writing the real and imaginary parts of $\widetilde{C}_{\varphi_{a,1}}(z)$ in terms of polar coordinates we have,
\begin{align} 
\mathrm{Re}(\widetilde{C}_{\varphi_{a,1}})(r,\theta)=&\left(\dfrac{1-r^2}{1-2ar\cos\theta+r^2} \right)^2\left(\left(1-ar\cos\theta \right)^2-\left(ar\sin\theta\right)^2 \right) \label{polarform}  \\
\mathrm{Im}(\widetilde{C}_{\varphi_{a,1}})(r,\theta)=&2\left(\dfrac{1-r^2}{1-2ar\cos\theta+r^2} \right)^2\left(1-ar\cos\theta \right)\left(-ar\sin\theta\right). \nonumber
\end{align}
Suppose $\mathrm{Re}(\widetilde{C}_{\varphi_{a,1}})(r,\theta)=0$, then this is equivalent to $\left(1-ar\cos\theta \right)^2-\left(ar\sin\theta\right)^2 =0$. Further, 
\begin{align}
 \left(1-ar\cos\theta \right)^2-\left(ar\sin\theta\right)^2=0 \quad 
 \Leftrightarrow& \quad 1-2ar\cos\theta+a^2r^2(2\cos^2\theta-1)=0 \nonumber \\
 \Leftrightarrow& \quad 2a^2r^2\cos^2\theta-2ar\cos\theta+1-a^2r^2=0 \nonumber \\
 \Leftrightarrow& \quad \cos^2\theta-\dfrac{1}{ar}\cos\theta+\dfrac{1-a^2r^2}{2a^2r^2}=0. \label{quadraticequation}
\end{align}
The quadratic equation in terms of $\cos \theta$ has a solution for $ \frac{\sqrt{2}}{2r}\leq a.$ By our assumption that $\frac{\sqrt{2}}{2}<a<1$, we have $\frac{\sqrt{2}}{2}<r<1$. We will show that equation \eqref{quadraticequation} has a solution for all $\theta \in [0,2\pi]$ where $\theta$ depends on $a$ and $r$ satisfying the inequality $\frac{\sqrt{2}}{2r} \leq a < 1.$ Solving for the solution of $\cos \theta$, we have 
\begin{align*}
\cos \theta = \dfrac{1 \pm \sqrt{2a^2r^2-1}}{2ar} \Rightarrow& -1 \leq \dfrac{1 \pm \sqrt{2a^2r^2-1}}{2ar} \leq 1 \\
\Rightarrow& -2ar-1 \leq \pm \sqrt{2a^2r^2-1} \leq 2ar-1.
\end{align*}
This indicates that both inequalities  
\begin{align}\label{inequality1}
-2ar-1 \leq \sqrt{2a^2r^2-1} \leq 2ar-1, \text{and}    
\end{align}
\begin{align}\label{inequality2}
-2ar-1 \leq - \sqrt{2a^2r^2-1} \leq 2ar-1    
\end{align}
must hold true.
\noindent In inequality \eqref{inequality1}, $-2ar-1 \leq \sqrt{2a^2r^2-1}$ holds true. Since $a^2r^2-2ar+1=(ar-1)^2 \geq 0$, we have
\begin{align*}
0 \leq a^2r^2-2ar+1  \Leftrightarrow&\  0 \leq 2a^2r^2-4ar+2  \\
\Leftrightarrow&\ 2a^2r^2-1 \leq 4a^2r^2-4ar+1  \\
\Leftrightarrow&\ 2a^2r^2-1 \leq (2ar-1)^2 \\
\Leftrightarrow&\ \sqrt{2a^2r^2-1} \leq 2ar-1. \\
\end{align*}
In inequality \eqref{inequality2}, since $2ar-1 \geq 0$ we have $- \sqrt{2a^2r^2-1} \leq 2ar-1$. Since $a^2r^2+2ar+1=(ar+1)^2 \geq 0$, we have
\begin{align*}
0 \leq a^2r^2+2ar+1 \Leftrightarrow&\ 0 \leq 2a^2r^2+4ar+2 \\
\Leftrightarrow&\ 2a^2r^2-1 \leq 4a^2r^2+4ar+1  \\
\Leftrightarrow&\ 2a^2r^2-1 \leq (-2ar-1)^2  \\
\Leftrightarrow&\ -\sqrt{2a^2r^2-1} \leq -2ar-1. \\
\end{align*}
The affirmation of inequalities \eqref{inequality1} and \eqref{inequality2} shows that for $\frac{\sqrt{2}}{2r} \leq a < 1$ there exists a $\theta(a,r)$, that depends on $a$ and $r$, satisfying equation \eqref{quadraticequation}. Thus, we have $\mathrm{Re}(\widetilde{C}_{\varphi_{a,1}})(r,\theta(a,r))=0$ and $\mathrm{Im}(\widetilde{C}_{\varphi_{a,1}})(r,\theta(a,r)) \neq 0$, showing that
$\widetilde{C}_{{\varphi}_{a,1}}(r,\theta(a,r))\neq 0$ lies on the imaginary axis. By Corollary \ref{propconj}, $-\widetilde{C}_{{\varphi}_{a,1}}(r,\theta(a,r))$ is also in the Berezin range.  By the assumption that $B(C_{\varphi_{a,1}})$ is convex, there is a line segment connecting $\widetilde{C}_{{\varphi}_{a,1}}(r,\theta(a,r))$ and $-\widetilde{C}_{{\varphi}_{a,1}}(r,\theta(a,r))$ is contained in the Berezin range which indicates that $0$ is in the Berezin range as well. This is a contradiction since if $\widetilde{C}_{{\varphi}_{a,1}}(z)=0$ then $|z|=1$. Therefore, if $B(C_{\varphi_{a,1}})$ is convex, then $0 \leq a \leq \frac{\sqrt{2}}{2}$.

If $a=0$, then $B(C_{\varphi_{a,1}})=\{1\}$ is convex. Let $0< a \leq \frac{\sqrt{2}}{2}$, we want to show that $B(C_{\varphi_{a,1}})$ is convex. First, we want to show that for $0< r <1$, $B(C_{\varphi_{a,1}})|_{re^{i\theta}}=\gamma(\theta)$ where $\gamma$ is a smooth loop, $\gamma(0)=\gamma(2\pi)$ and $\gamma'(\theta) \neq 0$ exists. From \eqref{polarform}, 
\[\mathrm{Re}(\widetilde{C}_{\varphi_{a,1}})(r,0)=\mathrm{Re}(\widetilde{C}_{\varphi_{a,1}})(r,2\pi)=\left(\dfrac{1-r^2}{1-2ar+r^2}\right)^2\left(1-ar\right)^2\text{ and }\mathrm{Re}(\widetilde{C}_{\varphi_{a,1}})(r,\pi)=\left(\dfrac{1-r^2}{1+2ar+r^2}\right)^2\left(1+ar\right)^2,\]
where $\mathrm{Re}(\widetilde{C}_{\varphi_{a,1}})(r,0) > \mathrm{Re}(\widetilde{C}_{\varphi_{a,1}})(r,\pi)$ holds. For $\theta \in [0,2\pi]$,
\begin{align*}
\widetilde{C}_{\varphi_{a,1}}(0,\theta)=1\ \ \mathrm{and}\ \ \ \ \widetilde{C}_{\varphi_{a,1}}(1,\theta)=0.
\end{align*}
For $\theta \in (0,\pi)$, we observe that $\mathrm{Im}(\widetilde{C}_{\varphi_{a,1}})(r,\theta) < 0$. Evaluating $\dfrac{\partial}{\partial\theta}\mathrm{Re}(\widetilde{C}_{\varphi_{a,1}})$, we have
\begin{align}\label{DerReBR}
\dfrac{\partial}{\partial\theta}\mathrm{Re}(\widetilde{C}_{\varphi_{a,1}})(r,\theta)=&-\dfrac{4ar\sin(\theta)(r-1)^2(r+1)^2\left[-ar^3\cos(\theta)-\dfrac{1}{2}+\left(a^2+\dfrac{1}{2}\right)r^2\right]}{(2ar\cos(\theta)-r^2-1)^3}.
\end{align}
For a small $\varepsilon >0$, $\dfrac{\partial}{\partial\theta}\mathrm{Re}(\widetilde{C}_{\varphi_{a,1}})(r,\varepsilon) < 0$. By continuity of the Berezin range on $\theta \in [0,\pi]$, we have a curve $\gamma_1(\theta) \subset \{z \in \mathbb{C}|\ \mathrm{Im}(z) \leq 0\}$ where $\gamma_1(0)$ and $ \gamma_{1}(\pi)$ are real valued, and $\gamma_1(\pi) < \gamma_{1}(0)$. Since $-\mathrm{Im}(\widetilde{C}_{\varphi_{a,1}})(r,\theta)=\mathrm{Im}(\widetilde{C}_{\varphi_{a,1}})(r,-\theta)$ and $\mathrm{Re}(\widetilde{C}_{\varphi_{a,1}})(r,\theta)=\mathrm{Re}(\widetilde{C}_{\varphi_{a,1}})(r,-\theta)$ for $\theta \in [\pi,2\pi]$, the Berezin range traces a curve $\gamma_2 \subset \{z \in \mathbb{C}|\mathrm{Im}(z) \geq 0\}$ where $\gamma_2(\pi)$ and $ \gamma_{2}(2\pi)$ are real valued, and $\gamma_{2}(\pi) < \gamma_{2}(2\pi)$. Moreover, $\gamma_1$ and $\gamma_2$ are conjugate symmetric, that is $\overline{\gamma_1}=\gamma_2$. For any $r \in (0,1)$, we have $ B(C_{\varphi_{a,1}})|_{re^{i\theta}}=\gamma_{1,r}(\theta) \cup \gamma_{2,r}(\theta)=\gamma_r(\theta)$ which shows that the restriction of the Berezin range to circles on the disk is a loop.

Next, we want to show that the Berezin range does not contain holes. From \eqref{polarform}, we have 
\[\dfrac{\partial}{\partial r}\mathrm{Re}(\widetilde{C}_{\varphi_{a,1}})(r,0)=\dfrac{8(r+1)(ar-1)(r-1)\left(a^2r^3-\dfrac{1}{4}ar^4-\dfrac{3}{2}ar^2-\dfrac{1}{4}a+r\right)}{(2ar-r^2-1)^3}.\]
We then solve $\dfrac{\partial}{\partial r}\mathrm{Re}(\widetilde{C}_{\varphi_{a,1}})(0,0) =2a>0$, indicating that the Berezin range is increasing at $r=0$. Since
\begin{align*}
\widetilde{C}_{\varphi_{a,1}}(0,0)=1\ \ \mathrm{and}\ \ \ \ \widetilde{C}_{\varphi_{a,1}}(1,0)=0.
\end{align*}
and the Berezin range is increasing at $r=0$, by the extreme value theorem, there is an absolute maximum that the Berezin range achieves at $r_0$ for $0<r_0<1$. The loop traced by $B(C_{\varphi_{a,1}})|_{re^{i\theta}}=\gamma_r(\theta)$ starts at $1$ then traverses to the right until reaching $B(C_{\varphi_{a,1}})|_{r_0e^{i\theta}}=\gamma_{r_0}(\theta)$. After reaching the absolute maximum at $r_0$, the loops change direction traversing to the left going to zero. Assuming that the Berezin range has a hole impedes the loops from traversing $0$ to $1$ continuously. Thus, $B(C_{\varphi_{a,1}})$ has no holes.

Now, we focus at the boundary of the Berezin range, $b(B(C_{\varphi_{a,1}})) \subset \{z \in \mathbb{C}|\mathrm{Im}(\mathbb{C}) \leq 0\} $ and show that it is concave up. By complex conjugate symmetry, this shows that $b(B(C_{\varphi_{a,1}})) \subset \{z \in \mathbb{C}|\mathrm{Im}(\mathbb{C}) \geq 0\} $ is concave down. Let $r$ be fixed, we want to solve the extrema of $\mathrm{Im}(\widetilde{C}_{\varphi_{a,1}})(r,\theta)$. 
We observe that, 
\[\dfrac{\partial}{\partial\theta}\mathrm{Im}(\widetilde{C}_{\varphi_{a,1}})(r,\theta)=\dfrac{2ar[-2ar^3\cos^2(\theta)+(2a^2r^2+r^2+1)\cos(\theta)+ar^3-3ar](r-1)^2(r+1)^2}{(2ar\cos(\theta)-r^2-1)^3}.\]
We take note that $-2ar^3\cos^2(\theta)+(2a^2r^2+r^2+1)\cos(\theta)+ar^3-3ar$ is quadratic in terms of $\cos(\theta)$ and when equated to zero, has a solution $\theta_1 \in (0,\pi)$ . Specifically, the solution $\theta_1$ is a critical value  dependent to $r$ such that $\mathrm{Im}(\widetilde{C}_{\varphi_{a,1}})(r,\theta_1)$ achieves a minimum. Using the quadratic formula, one can compute that 
\begin{align}\label{theta_1}
\theta_1 = \pi - \cos^{-1}\left(\dfrac{-2a^2r^2-r^2-1+\sqrt{4a^4r^4+8a^2r^6-20r^4a^2+4a^2r^2+r^4+2r^2+1}}{4ar^3}\right)
\end{align}
Since $\mathrm{Im}(\widetilde{C}_{\varphi_{a,1}})(r,0)=\mathrm{Im}(\widetilde{C}_{\varphi_{a,1}})(r,\pi)=0$ for $0<r<1$, we achieve global minimum at $\theta_1$.
Consider  $\min(\widetilde{C}_{\varphi_{a,1}}): [0,1] \rightarrow \mathbb{R}^2$ defined by 
\[\min(\widetilde{C}_{\varphi_{a,1}})(r)=(\mathrm{Re}(\widetilde{C}_{\varphi_{a,1}})(r,\theta_1),\mathrm{Im}(\widetilde{C}_{\varphi_{a,1}})(r,\theta_1)).\]
Since $\mathrm{Im}(\widetilde{C}_{\varphi_{a,1}})(r,\theta_1)$ achieves a minimum for any $r \in (0,1)$, 
\[\{\min(\widetilde{C}_{\varphi_{a,1}})(r)|\ r \in [0,1]\} = b(B(C_{\varphi_{a,1}})) \cap \{z \in \mathbb{C}|\ \mathrm{Im}(z) \leq 0\} \]
We then verify concavity using second derivative test, 
\[\dfrac{d^2\mathrm{Im}(\widetilde{C}_{\varphi_{a,1}})}{d\mathrm{Re}(\widetilde{C}_{\varphi_{a,1}})^2}=\dfrac{\frac{d}{dr}\left(\frac{d\mathrm{Im}(\widetilde{C}_{\varphi_{a,1}})/dr}{d\mathrm{Re}(\widetilde{C}_{\varphi_{a,1}})/dr}\right)}{\frac{dRe(\widetilde{C}_{\varphi_{a,1}})}{dr}} > 0.\]
Thus, the parametric curve is concave up and by the fact that $\min(\widetilde{C}_{\varphi_{a,1}})(0)=(1,0)$ and $\min(\widetilde{C}_{\varphi_{a,1}})(1)=(0,0)$, there exists $r=r_{\mathrm{min}}$ where $\mathrm{Im}(\widetilde{C}_{\varphi_{a,1}})(r_{\mathrm{min}},\theta_1)$ is the global minimum. By complex conjugate symmetry the boundary curve $b(B(C_{\varphi_{a,1}})) \cap \{z \in \mathbb{C}|\ \mathrm{Im}(z) \geq 0\}$ is concave down and $\mathrm{Im}(\widetilde{C}_{\varphi_{a,1}})$ achieves a global maximum at $r=r_{\mathrm{min}}$.

Thus, we have shown that the loops start at $1$, and as it travels to the right and travel back to the left until it reaches the origin, the peak of the loops increases in magnitude and reaches an extremum before it decreases to zero. The peaks of the loops trace a subboundary of the Berezin range which is convex. Another part of the boundary of the Berezin range is due to the loop $B(C_{\varphi_{a,1}})|_{r_1e^{i\theta}}$ where $r_1$ is the radius so that $\mathrm{Re}(\widetilde{C}_{\varphi_{a,1}})(r_1,0)$ is the global maximum. We want to observe the contour defined by $\widetilde{C}_{\varphi_{a,1}}(r_1,\theta): [0,\theta_1] \rightarrow \mathbb{R}^2 $ where $\theta_1$ is defined in \eqref{theta_1}.  Using second derivative test to verify concavity, 
\[\dfrac{d^2\mathrm{Im}(\widetilde{C}_{\varphi_{a,1}})}{d\mathrm{Re}(\widetilde{C}_{\varphi_{a,1}})^2}=\dfrac{\frac{d}{d\theta}\left(\frac{d\mathrm{Im}(\widetilde{C}_{\varphi_{a,1}})/d\theta}{d\mathrm{Re}(\widetilde{C}_{\varphi_{a,1}})/d\theta}\right)}{\frac{dRe(\widetilde{C}_{\varphi_{a,1}})}{d\theta}} > 0.\]
This shows that the parametric curve is concave up, and by complex conjugate symmetry, the curve above is concave down. With all of this information, we have shown that the Berezin range $B(C_{\varphi_{a,1}})$ is convex. By Proposition \ref{auto_on_disk}, $B(C_{\varphi_{a,1}})$ is convex for $0< |a| \leq \frac{\sqrt{2}}{2}$.  
\end{proof}

\end{document}
